\documentclass[11pt,a4paper]{ctexart}
\usepackage[margin=2.2cm]{geometry}
\usepackage{xcolor}
\usepackage{enumitem}
\usepackage{amssymb}
\usepackage[normalem]{ulem}
\usepackage{titlesec}
\usepackage{hyperref}
\definecolor{navy}{HTML}{1A365D}
\definecolor{ink}{HTML}{1F2933}
\definecolor{muted}{HTML}{66788A}
\definecolor{soft}{HTML}{EEF4FA}
\definecolor{gold}{HTML}{8A5A12}
\hypersetup{colorlinks=true,linkcolor=navy,urlcolor=navy}
\setlist{nosep}
\titleformat{\section}{\Large\bfseries\color{navy}}{}{0pt}{}
\titleformat{\subsection}{\large\bfseries\color{navy}}{}{0pt}{}
\pagestyle{plain}

\title{\textbf{八上不定代词专项}\\[4pt]
\large 人教版（新课纲版）八年级上册 · 不定代词（Indefinite Pronouns）}
\author{初中英语资料站}
\date{2026-11-12}
\begin{document}
\maketitle

\section*{使用说明}
\begin{itemize}[leftmargin=1.4em]
\item 本讲义对应人教版（新课纲版）八年级上册语法要求：不定代词的用法。
\item 共 4 讲语法、6 级 48 题。
\item 在线版：\url{https://1901151458-prog.github.io/english-materials/g8a-pronouns/}
\end{itemize}

\section{some 和 any}
some 多用于肯定句，any 多用于否定句、疑问句和条件从句。两者既可代替或修饰可数名词，也可代替或修饰不可数名词。
\begin{itemize}[leftmargin=1.4em]
\item A: Are there \underline{any} apples in the fridge? B: Yes, there are \underline{some}. / No, there aren't \underline{any}.
\item A: Is there \underline{any} water in the bottle? B: Yes, there is \underline{some} water.
\end{itemize}
\textbf{例外：}在「Would you like some tea?」这类疑问句中用 \textbf{some}，因为说话人希望得到对方\textbf{肯定的答复}。

\section{many 和 much}
many 只代替或修饰\textbf{可数名词复数}，much 只代替或修饰\textbf{不可数名词}。
\begin{itemize}[leftmargin=1.4em]
\item In the future, \underline{many} people will live in outer space.
\item There isn't \underline{much} water left in the bottle.
\item I'm afraid the kids spend too \underline{much} time online.
\end{itemize}

\section{few / a few / little / a little}
\begin{center}
\begin{tabular}{p{2.6cm}p{2.8cm}p{3.6cm}p{4.4cm}}
\hline
\textbf{词} & \textbf{修饰} & \textbf{含义} & \textbf{例句}\\
\hline
few & 可数复数 & 几乎没有（否定） & \underline{Few} people know what happened.\\
a few & 可数复数 & 一些（肯定） & We stayed there for \underline{a few} days.\\
little & 不可数 & 几乎没有（否定） & There is \underline{little} hope of success.\\
a little & 不可数 & 一些（肯定） & There is \underline{a little} time left.\\
\hline
\end{tabular}
\end{center}
\textbf{关键：}没有 a 的 few / little 侧重\textbf{否定}；有 a 的 a few / a little 侧重\textbf{肯定}。

\section{复合不定代词}
由 some、any、no、every 加上 -body、-one、-thing 构成。除 no one 外都写成一个词。复合不定代词只相当于名词，用作主语、宾语或表语；\textbf{作主语时都作单数看待}。
\begin{itemize}[leftmargin=1.4em]
\item \underline{Everyone} says Yulong Mountain is beautiful.
\item \underline{Nobody} knows what the future will be like.
\item I'm on vacation in Hangzhou and \underline{everything} is fantastic.
\end{itemize}
\textbf{定语后置：}复合不定代词如果有定语修饰，定语须置于其后。
\begin{itemize}[leftmargin=1.4em]
\item Did you do \underline{anything fun} on your holiday?
\item I did \underline{nothing special}.
\item Yesterday, I met \underline{someone interesting}.
\end{itemize}


\newpage
\section{分级练习（48 题）}
\noindent 按难度分级排列，先做基础，再做成分与扩展，最后做综合题。

\subsection*{第 1 级 · some 和 any}
\noindent\textcolor{gray}{some 多用于肯定句，any 多用于否定句、疑问句和条件从句。}\par\smallskip
\begin{enumerate}[leftmargin=*,itemsep=4pt]
\item Would you like \_\_\_ tea?
\begin{enumerate}[label=\Alph*.,leftmargin=2em,itemsep=1pt,topsep=2pt]
\item some
\item any
\item much
\item many
\end{enumerate}
\item I don’t have \_\_\_ money with me.
\begin{enumerate}[label=\Alph*.,leftmargin=2em,itemsep=1pt,topsep=2pt]
\item some
\item any
\item many
\item much
\end{enumerate}
\item Are there \_\_\_ books on the desk?
\begin{enumerate}[label=\Alph*.,leftmargin=2em,itemsep=1pt,topsep=2pt]
\item some
\item any
\item many
\item much
\end{enumerate}
\item Could I have \_\_\_ water, please?
\begin{enumerate}[label=\Alph*.,leftmargin=2em,itemsep=1pt,topsep=2pt]
\item some
\item any
\item many
\item much
\end{enumerate}
\item He didn’t buy \_\_\_ fruit yesterday.
\begin{enumerate}[label=\Alph*.,leftmargin=2em,itemsep=1pt,topsep=2pt]
\item some
\item any
\item many
\item much
\end{enumerate}
\item Do you have \_\_\_ questions?
\begin{enumerate}[label=\Alph*.,leftmargin=2em,itemsep=1pt,topsep=2pt]
\item some
\item any
\item many
\item much
\end{enumerate}
\item If you have \_\_\_ problems, tell me.
\begin{enumerate}[label=\Alph*.,leftmargin=2em,itemsep=1pt,topsep=2pt]
\item some
\item any
\item many
\item much
\end{enumerate}
\item A: Is there \_\_\_ water in the bottle? B: Yes, there is \_\_\_.
\begin{enumerate}[label=\Alph*.,leftmargin=2em,itemsep=1pt,topsep=2pt]
\item any; some
\item some; any
\item any; any
\item some; some
\end{enumerate}
\end{enumerate}

\subsection*{第 2 级 · many 和 much}
\noindent\textcolor{gray}{many 修饰可数名词复数，much 修饰不可数名词。}\par\smallskip
\begin{enumerate}[leftmargin=*,itemsep=4pt]
\item In the future, \_\_\_ people will live in outer space.
\begin{enumerate}[label=\Alph*.,leftmargin=2em,itemsep=1pt,topsep=2pt]
\item many
\item much
\item few
\item little
\end{enumerate}
\item There isn’t \_\_\_ water left in the bottle.
\begin{enumerate}[label=\Alph*.,leftmargin=2em,itemsep=1pt,topsep=2pt]
\item many
\item much
\item few
\item little
\end{enumerate}
\item I’m afraid the kids spend too \_\_\_ time online.
\begin{enumerate}[label=\Alph*.,leftmargin=2em,itemsep=1pt,topsep=2pt]
\item many
\item much
\item few
\item little
\end{enumerate}
\item How \_\_\_ students are there in your class?
\begin{enumerate}[label=\Alph*.,leftmargin=2em,itemsep=1pt,topsep=2pt]
\item many
\item much
\item few
\item little
\end{enumerate}
\item How \_\_\_ milk do you need?
\begin{enumerate}[label=\Alph*.,leftmargin=2em,itemsep=1pt,topsep=2pt]
\item many
\item much
\item few
\item little
\end{enumerate}
\item She has \_\_\_ friends in this city.
\begin{enumerate}[label=\Alph*.,leftmargin=2em,itemsep=1pt,topsep=2pt]
\item many
\item much
\item few
\item little
\end{enumerate}
\item We don’t have \_\_\_ time before the train leaves.
\begin{enumerate}[label=\Alph*.,leftmargin=2em,itemsep=1pt,topsep=2pt]
\item many
\item much
\item few
\item little
\end{enumerate}
\item There are \_\_\_ books in the school library.
\begin{enumerate}[label=\Alph*.,leftmargin=2em,itemsep=1pt,topsep=2pt]
\item many
\item much
\item few
\item little
\end{enumerate}
\end{enumerate}

\subsection*{第 3 级 · few / a few / little / a little}
\noindent\textcolor{gray}{few、little 侧重否定；a few、a little 侧重肯定。}\par\smallskip
\begin{enumerate}[leftmargin=*,itemsep=4pt]
\item \_\_\_ people know what happened. It’s a secret.
\begin{enumerate}[label=\Alph*.,leftmargin=2em,itemsep=1pt,topsep=2pt]
\item Few
\item A few
\item Little
\item A little
\end{enumerate}
\item We stayed in the old town for \_\_\_ days.
\begin{enumerate}[label=\Alph*.,leftmargin=2em,itemsep=1pt,topsep=2pt]
\item few
\item a few
\item little
\item a little
\end{enumerate}
\item Don’t worry. There is \_\_\_ time left.
\begin{enumerate}[label=\Alph*.,leftmargin=2em,itemsep=1pt,topsep=2pt]
\item few
\item a few
\item little
\item a little
\end{enumerate}
\item There is now \_\_\_ hope of success.
\begin{enumerate}[label=\Alph*.,leftmargin=2em,itemsep=1pt,topsep=2pt]
\item few
\item a few
\item little
\item a little
\end{enumerate}
\item He has \_\_\_ money, so he can’t buy the bike.
\begin{enumerate}[label=\Alph*.,leftmargin=2em,itemsep=1pt,topsep=2pt]
\item few
\item a few
\item little
\item a little
\end{enumerate}
\item I have \_\_\_ questions to ask you.
\begin{enumerate}[label=\Alph*.,leftmargin=2em,itemsep=1pt,topsep=2pt]
\item few
\item a few
\item little
\item a little
\end{enumerate}
\item \_\_\_ students were late today — only two.
\begin{enumerate}[label=\Alph*.,leftmargin=2em,itemsep=1pt,topsep=2pt]
\item Few
\item A few
\item Little
\item A little
\end{enumerate}
\item Can I have \_\_\_ sugar in my coffee?
\begin{enumerate}[label=\Alph*.,leftmargin=2em,itemsep=1pt,topsep=2pt]
\item few
\item a few
\item little
\item a little
\end{enumerate}
\end{enumerate}

\subsection*{第 4 级 · 复合不定代词}
\noindent\textcolor{gray}{some/any/no/every + -body/-one/-thing，作主语用单数。}\par\smallskip
\begin{enumerate}[leftmargin=*,itemsep=4pt]
\item \_\_\_ says Yulong Mountain is beautiful.
\begin{enumerate}[label=\Alph*.,leftmargin=2em,itemsep=1pt,topsep=2pt]
\item Everyone
\item Every one
\item Every
\item Each
\end{enumerate}
\item \_\_\_ knows what the future will be like.
\begin{enumerate}[label=\Alph*.,leftmargin=2em,itemsep=1pt,topsep=2pt]
\item Nobody
\item None
\item No one of
\item Nothing
\end{enumerate}
\item I’m on vacation in Hangzhou and \_\_\_ is fantastic.
\begin{enumerate}[label=\Alph*.,leftmargin=2em,itemsep=1pt,topsep=2pt]
\item everything
\item something
\item nothing
\item anything
\end{enumerate}
\item Is there \_\_\_ in the box?
\begin{enumerate}[label=\Alph*.,leftmargin=2em,itemsep=1pt,topsep=2pt]
\item anything
\item something
\item nothing
\item everything
\end{enumerate}
\item \_\_\_ is impossible if you try your best.
\begin{enumerate}[label=\Alph*.,leftmargin=2em,itemsep=1pt,topsep=2pt]
\item Nothing
\item Something
\item Anything
\item Everything
\end{enumerate}
\item Is there \_\_\_ interesting in today’s newspaper?
\begin{enumerate}[label=\Alph*.,leftmargin=2em,itemsep=1pt,topsep=2pt]
\item anything
\item something
\item nothing
\item everything
\end{enumerate}
\item \_\_\_ of the students came to school because of the heavy snow.
\begin{enumerate}[label=\Alph*.,leftmargin=2em,itemsep=1pt,topsep=2pt]
\item None
\item Nobody
\item No one
\item Nothing
\end{enumerate}
\item There is \_\_\_ wrong with my bike. It works well.
\begin{enumerate}[label=\Alph*.,leftmargin=2em,itemsep=1pt,topsep=2pt]
\item nothing
\item something
\item anything
\item everything
\end{enumerate}
\end{enumerate}

\subsection*{第 5 级 · 定语后置}
\noindent\textcolor{gray}{复合不定代词有定语修饰时，定语放在它后面。}\par\smallskip
\begin{enumerate}[leftmargin=*,itemsep=4pt]
\item Did you do anything \_\_\_ on your holiday?
\begin{enumerate}[label=\Alph*.,leftmargin=2em,itemsep=1pt,topsep=2pt]
\item fun
\item funny
\item funly
\item funs
\end{enumerate}
\item I did nothing \_\_\_.
\begin{enumerate}[label=\Alph*.,leftmargin=2em,itemsep=1pt,topsep=2pt]
\item special
\item specially
\item specialness
\item specials
\end{enumerate}
\item Yesterday, I met someone \_\_\_.
\begin{enumerate}[label=\Alph*.,leftmargin=2em,itemsep=1pt,topsep=2pt]
\item interesting
\item interested
\item interest
\item interestingly
\end{enumerate}
\item Is there anything \_\_\_ in the newspaper?
\begin{enumerate}[label=\Alph*.,leftmargin=2em,itemsep=1pt,topsep=2pt]
\item new
\item newly
\item news
\item newer
\end{enumerate}
\item She wants something \_\_\_ to eat.
\begin{enumerate}[label=\Alph*.,leftmargin=2em,itemsep=1pt,topsep=2pt]
\item sweet
\item sweetly
\item sweetness
\item sweets
\end{enumerate}
\item Did you see anybody \_\_\_ at the party?
\begin{enumerate}[label=\Alph*.,leftmargin=2em,itemsep=1pt,topsep=2pt]
\item famous
\item famously
\item fame
\item famousness
\end{enumerate}
\item There is nothing \_\_\_ with your computer.
\begin{enumerate}[label=\Alph*.,leftmargin=2em,itemsep=1pt,topsep=2pt]
\item wrong
\item wrongly
\item wrongness
\item wrongs
\end{enumerate}
\item I have something \_\_\_ to tell you.
\begin{enumerate}[label=\Alph*.,leftmargin=2em,itemsep=1pt,topsep=2pt]
\item important
\item importantly
\item importance
\item importants
\end{enumerate}
\end{enumerate}

\subsection*{第 6 级 · 综合提升}
\noindent\textcolor{gray}{some/any、few/little、复合不定代词与 either/neither/none 混辨。}\par\smallskip
\begin{enumerate}[leftmargin=*,itemsep=4pt]
\item Would you like \_\_\_ to drink?
\begin{enumerate}[label=\Alph*.,leftmargin=2em,itemsep=1pt,topsep=2pt]
\item something
\item anything
\item nothing
\item everything
\end{enumerate}
\item I don’t have \_\_\_ to say about it.
\begin{enumerate}[label=\Alph*.,leftmargin=2em,itemsep=1pt,topsep=2pt]
\item something
\item anything
\item nothing
\item everything
\end{enumerate}
\item \_\_\_ of the two answers is correct. Both are wrong.
\begin{enumerate}[label=\Alph*.,leftmargin=2em,itemsep=1pt,topsep=2pt]
\item Neither
\item Either
\item None
\item Both
\end{enumerate}
\item There are trees on \_\_\_ side of the street.
\begin{enumerate}[label=\Alph*.,leftmargin=2em,itemsep=1pt,topsep=2pt]
\item either
\item both
\item all
\item none
\end{enumerate}
\item \_\_\_ of the three students passed the exam.
\begin{enumerate}[label=\Alph*.,leftmargin=2em,itemsep=1pt,topsep=2pt]
\item None
\item Neither
\item Either
\item Both
\end{enumerate}
\item \_\_\_ of us has a dictionary.
\begin{enumerate}[label=\Alph*.,leftmargin=2em,itemsep=1pt,topsep=2pt]
\item Each
\item Every
\item Both
\item All
\end{enumerate}
\item \_\_\_ the students are here today.
\begin{enumerate}[label=\Alph*.,leftmargin=2em,itemsep=1pt,topsep=2pt]
\item All
\item Both
\item Each
\item Every
\end{enumerate}
\item A: Are there any apples in the fridge? B: No, there aren’t \_\_\_.
\begin{enumerate}[label=\Alph*.,leftmargin=2em,itemsep=1pt,topsep=2pt]
\item any
\item some
\item many
\item much
\end{enumerate}
\end{enumerate}

\newpage
\section*{答案与解析}
\small
\begin{enumerate}[leftmargin=1.6em,itemsep=2pt]
\item \textbf{U1-1 A}：希望对方肯定答复时用 some，不用 any；教材例句。
\item \textbf{U1-2 B}：否定句用 any。
\item \textbf{U1-3 B}：疑问句用 any。
\item \textbf{U1-4 A}：请求并希望得到肯定答复，用 some。
\item \textbf{U1-5 B}：否定句用 any。
\item \textbf{U1-6 B}：疑问句用 any。
\item \textbf{U1-7 B}：条件从句中用 any；教材说明 any 也用于条件从句。
\item \textbf{U1-8 A}：疑问句用 any，肯定回答用 some；教材例句。
\item \textbf{U2-1 A}：people 可数复数，用 many；教材例句。
\item \textbf{U2-2 B}：water 不可数，用 much；教材例句。
\item \textbf{U2-3 B}：time 不可数，用 much；教材例句。
\item \textbf{U2-4 A}：students 可数复数，用 many。
\item \textbf{U2-5 B}：milk 不可数，用 much。
\item \textbf{U2-6 A}：friends 可数复数，用 many。
\item \textbf{U2-7 B}：time 不可数，用 much。
\item \textbf{U2-8 A}：books 可数复数，用 many。
\item \textbf{U3-1 A}：people 可数 + 否定含义（几乎没有），用 Few；教材例句。
\item \textbf{U3-2 B}：days 可数 + 肯定含义（一些），用 a few；教材例句。
\item \textbf{U3-3 D}：time 不可数 + 肯定含义（还有一些），用 a little；教材例句。
\item \textbf{U3-4 C}：hope 不可数 + 否定含义（几乎没有），用 little；教材例句。
\item \textbf{U3-5 C}：money 不可数 + 否定含义，用 little。
\item \textbf{U3-6 B}：questions 可数复数 + 肯定含义，用 a few。
\item \textbf{U3-7 A}：可数 + 「几乎没有」，用 Few。
\item \textbf{U3-8 D}：sugar 不可数 + 「一点」，用 a little。
\item \textbf{U4-1 A}：复合不定代词作主语，谓语用单数；教材例句。
\item \textbf{U4-2 A}：nobody 作主语，谓语用单数；教材例句。
\item \textbf{U4-3 A}：表示「一切」，用 everything；教材例句。
\item \textbf{U4-4 A}：疑问句用 anything。
\item \textbf{U4-5 A}：Nothing is impossible「没有不可能」，用 nothing。
\item \textbf{U4-6 A}：疑问句用 anything。
\item \textbf{U4-7 A}：none of + 名词，表示「都不」；nobody / no one 不能接 of。
\item \textbf{U4-8 A}：There is nothing wrong with ...「没有毛病」。
\item \textbf{U5-1 A}：修饰复合不定代词的定语后置：anything fun；教材例句。
\item \textbf{U5-2 A}：nothing special「没什么特别的」，定语后置；教材例句。
\item \textbf{U5-3 A}：someone interesting，定语后置；教材例句。
\item \textbf{U5-4 A}：anything new，定语后置。
\item \textbf{U5-5 A}：something sweet，定语后置。
\item \textbf{U5-6 A}：anybody famous，定语后置。
\item \textbf{U5-7 A}：nothing wrong，定语后置。
\item \textbf{U5-8 A}：something important，定语后置。
\item \textbf{U6-1 A}：期待肯定答复的疑问句用 something。
\item \textbf{U6-2 B}：否定句用 anything。
\item \textbf{U6-3 A}：两者都不，用 neither；教材列出「neither 没人或物（指两个当中）」。
\item \textbf{U6-4 A}：either「任何一个（指两个当中）」，后接单数 side。
\item \textbf{U6-5 A}：三者以上都不，用 none；教材列出「none 没人或物（指两个以上）」。
\item \textbf{U6-6 A}：each of + 名词，谓语用单数。
\item \textbf{U6-7 A}：all the students 复数，谓语用 are。
\item \textbf{U6-8 A}：否定答句用 any；教材例句。
\end{enumerate}

\end{document}
